\chapter{基于的青光眼手术视频导航}

\section{引言}
青光眼（Glaucoma）作为全球范围内导致不可逆性失明的首要病因，
已成为严重的公共卫生负担\cite{flaxman2017global,sabanayagam2017global}。
微创青光眼手术凭借其创伤小、术后恢复快等优势，
已成为当前临床治疗的重要手段\cite{song2022minimally,zhang2023influence}。
然而，在眼科手术领域，手术空间极度狭小且解剖结构复杂。
前房角区域包含多层半透明且细微的生理结构，结构错综复杂，不同个体形态差异大，
即视为非青光眼的眼科专家，也同样难以快速识别\cite{dietlein2000morphological}。
统计数据显示，培养一名熟练的眼科医生需要漫长的学徒制教学周期，
而我国每 10 万人中眼科医生仅有不到4名，比例远低于发达国家\cite{nhc2022eyehealth, resnikoff2020estimated}，
优质医疗资源的紧缺与陡峭的手术学习曲线使得医疗事故风险依然存在。

在手术操作过程中，外科医生需要在复杂的动态场景下做出瞬时决策。
任何细微的识别延迟或步骤误判，都可能引发前房积血或眼压波动等严重并发症。
因此，实现在线手术阶段识别具有至关重要的临床意义。
在线手术阶段识别是构建“术中主动安全预警系统”的核心，
它不仅能为初级医生提供实时的规程指引，帮助其建立正确的认知框架，
还能作为下游任务（如风险预测、自动器械追踪）的基础，
从而真正将人工智能技术整合进手术流程，提升操作的标准化程度。

% 尽管针对普通外科等大尺度动作手术（如胸腹腔手术）的识别研究已取得一定进展
% \cite{endonet, czempiel2020tecno, wang2022autolaparo, lavanchy2024challenges, liu2025lovit}，
% 但
在微创青光眼手术这一高度精细且实时性要求极高的场景下，实现鲁棒的在线识别仍面临以下核心挑战：
\begin{enumerate}
    \item \textbf{时序分布失衡，短时动作难以召回}：
    手术视频具有显著的“长尾”数据特征，其流程通常由漫长的标准操作（如房角观察）和极其短暂的切换或精准交互操作组成。
    在这种分布极度不平衡的状态下，一些长阶段占据了数分钟时长，而决定阶段切换的关键动作（如注射粘弹剂）往往只有数秒。
    这导致模型极易产生预测偏向，倾向于维持高频的长阶段预测，
    而忽略低频但关键的短动作，导致关键切换节点的感知缺失，严重影响系统对手术状态判别的灵敏度。
    
    \item \textbf{外观高度相似，不同阶段难以区分}：
    在显微手术中，不同阶段的操作往往共用相同的手术器械，且背景环境几乎恒定（如始终聚焦于眼球且组织不会发生显著形变）。
    例如，“液体清洗”与“粘弹剂注入”在视觉特征上可能仅表现为细微的流体质感差异或者器械深入长短的差距。
    在在线识别场景下，模型无法像“事后复盘”那样查看后续画面来反向印证当前动作，只能依赖已有的画面信息。
    当不同阶段的视觉表现高度雷同，仅凭画面的时序特征很难实现精确区分，极易产生判断震荡。
\end{enumerate}


%% 现有工作不足：
% 现有的在线手术阶段识别研究主要侧重于构建复杂的时序模型（如 RNN、TCN 或 Transformer）以捕捉视频的长程依赖。
% 为满足实时性需求并解决数据分布问题，近期研究如 SKiT\cite{liu2023skit} 引入了高效的键聚类池化操作，
% 而 Meta-SurDiff\cite{li2025metasurdiff} 尝试利用分类扩散模型对预测分布进行细化。
% 尽管这些方法在处理通用手术流程方面取得了进展，但仍存在一定局限。

% 首先，针对“短时瞬态动作”的感知机制仍停留于静态语义层面。
% 现有的生成式或池化类方法本质上更倾向于对视频序列进行“全局分布拟合”，
% 即通过统计图像中器械与解剖结构的整体响应频率来定性，而忽略了帧间微弱的动态演化特征。
% 在青光眼手术中，诸如针尖刺入瞬间的特征突变往往被长时段的静态背景信息所淹没，
% 导致模型对关键动作切换点的感知缺乏物理意义上的“过程灵敏度”，难以在视觉干扰下实现鲁棒的特征增强。


%% 上下文切换管理：ProTAS（shen2024progress）（采用了Task Graph Learner），
%% SurgPlan++（chen2025surgplan++）（采用了FSM），ActionSwitch（注意针对overlap的设计）% 类似方案：ding2024neural（fsm）

% 其次，针对“视觉歧义”的手术规程约束缺乏环境适应性。
% 虽然 ProTAS\cite{shen2024protas} 与 NFSM\cite{ding2025nfsm} 分别尝试引入任务图与状态转移先验来强制维护时序连贯性，
% 但其机理主要建立在对训练集阶段跳转频率的“统计关联”建模上。
% 这种建模方式将规程逻辑视为一种硬性的后期修正，缺乏对实时场景中“环境置信度”与“内容逻辑”的深度解耦。
% 具体而言，模型无法区分当前的预测不确定性究竟是源于视频画面本身的质量降级（如术中出血或遮挡导致的感知失效），
% 还是来自动作序列与规程先验的逻辑冲突。
% 这种局限性导致模型在视觉清晰时可能受到逻辑先验的过度干扰，而在视觉模糊时又无法产生有效的“逻辑拉回”作用，
% 难以实现噪声外观特征与结构化规程知识之间的深度语义对齐。

% 综上所述，如何设计一种能感知瞬时动作动态差异的特征提取算子，
% 并构建一套能根据环境置信度自适应调整约束强度的逻辑校验引擎，是实现鲁棒在线识别的关键。


为解决上述挑战，本章在现有在线检测框架的基础上，提出了一种结合特征增强与逻辑修正的方法。
针对瞬时动作捕捉难题，我们设计了\textbf{自适应多尺度差分模块（A-MSDD）}。
该模块通过在线计算不同时距的特征差分，显式强化了瞬态变化信号，使模型能够精确锁定实时流中的阶段切换点。

针对语义混淆问题，本章进一步提出了\textbf{带门控的图逻辑插件（L-GBP）}。
该插件将手术逻辑转化为可进化的图表征，并配合双门控机制。
通过对当前环境置信度的实时评估，模型能自适应地利用逻辑先验（模拟医生经验）对视觉判别进行纠偏，显著增强了在线识别的一致性。

此外，本章提出了一种\textbf{自适应时序一致性损失函数（SA-Loss）}。
该损失函数能根据特征波动率动态调节监督力度，在保证在线识别低延迟的同时，有效平衡了切换边界的响应速度与阶段内部的平滑度。
通过上述模块的有机结合，本章构建了一个轻量、高效的在线手术阶段识别框架。


\section{问题定义}

在线动作检测（OAD）任务旨在从实时视频流中识别当前时刻正在发生的动作，并在每一帧即时预测其对应的类别标签。
给定一个持续输入的视频流 $V=\{f_t\}_{t=1}^T$，其中 $f_t$ 表示在时刻 $t$ 到达的视频帧，
任务的目标是根据历史观察及当前帧 $f_{1:t}$，在不依赖未来信息的前提下，
预测当前时刻 $t$ 的动作类别标签 $a_t \in \{0, 1, \dots, C\}$。在此定义中，
$C$ 表示预定义的动作类别总数，而标签 $0$ 通常被定义为“背景”类，
用于表示当前时刻没有任何目标动作发生。
不同于离线检测，在线动作检测要求模型在时间步 $t$ 实时输出预测结果 $\hat{a}_t$，即在每一个新的时间切片到达时，
建立从历史特征序列到当前动作状态的即时映射。

\section{本章方法}


\section{实验结果}

\subsection{数据集与评价指标}
\textbf{数据集构建}\quad

\textbf{评价指标}\quad

\textbf{实验细节}\quad

\subsection{性能比较}
% 与通用oad的性能比较（LSTR，TesTra, E2E-LOAD，CMerT，MerT）；
% 与在线手术识别模型的性能比较（SKViT，Meta-SurDiff，SurgPlan++）；


\subsection{消融分析}

\section{本章小结}